\documentclass[11pt,a4paper]{article}
\usepackage[margin=1in]{geometry}
\usepackage{times}
\usepackage{latexsym}
\usepackage{amsmath}
\usepackage{amssymb}
\usepackage{booktabs}
\usepackage{graphicx}
\usepackage{multirow}
\usepackage{xcolor}
\usepackage{url}
\usepackage{natbib}
\usepackage{enumitem}

\title{Principled Symbolic Consensus Improves LLM Reasoning}

\author{Anonymous}

\begin{document}
\maketitle

\begin{abstract}
Self-Consistency (SC) aggregates multiple LLM reasoning traces via majority voting but discards the logical structure within each trace.
When votes are tied, SC falls back to arbitrary heuristics (alphabetical ordering, random selection) with no principled resolution mechanism.
We introduce \textbf{SICA} (Symbolic Invariant Consensus Aggregation), which extracts first-order logic constraints from reasoning traces and resolves ties through weighted partial MAX-SAT solving.
On 804 questions spanning FOLIO and ProofWriter benchmarks, SICA achieves 71.3\% accuracy versus SC's 67.5\% (+3.8pp, $p=0.023$) with Qwen2.5-14B-Instruct.
SICA wins 75.9\% of tie cases, compared to 50\% expected from random tiebreaking.
However, we show that SICA's effectiveness is \emph{conditional on trace diversity}: with QwQ-32B, a reasoning model whose outputs are near-deterministic (multi-candidate rate $C_0=3.6\%$), SICA degenerates to SC.
We formalize this as the \textbf{operating regime} of symbolic consensus and demonstrate that $C_0$---a zero-cost diagnostic computable before running SICA---predicts the SICA--SC gap with Pearson $r=0.94$ across four model scales.
Our results establish when symbolic consensus adds value and when it does not, providing practitioners with a principled deployment criterion.
\end{abstract}

%-------------------------------------------------------------
\section{Introduction}
\label{sec:intro}

Large language models (LLMs) have demonstrated remarkable reasoning capabilities across diverse tasks, yet their outputs remain stochastic---sampling the same model multiple times on the same question yields different reasoning traces and sometimes different answers.
Self-Consistency (Wang et al., 2023) exploits this stochasticity by sampling $K$ traces and selecting the answer with the most votes (majority voting).
While effective, SC treats each reasoning trace as an opaque vote, discarding the rich logical structure contained within the chain-of-thought.

A fundamental limitation of majority voting emerges when votes are tied.
With even $K$ and questions that elicit diverse answers, ties occur in 8--15\% of cases.
SC resolves these ties through arbitrary heuristics---alphabetical ordering, random selection, or first-seen preference---none of which leverage the reasoning content.

We propose \textbf{SICA} (Symbolic Invariant Consensus Aggregation), a method that extracts first-order logic (FOL) constraints from reasoning traces and uses weighted partial MAX-SAT to find the answer most consistent with the extracted symbolic structure.
SICA operates as a drop-in replacement for SC's aggregation step: it uses the same $K$ sampled traces, adds a constraint extraction phase, and resolves both ties and close votes through principled symbolic optimization.

Our investigation reveals that SICA's effectiveness is not universal but is governed by a model property we call \textbf{trace diversity}, quantified by the multi-candidate rate $C_0$ (the fraction of questions where sampled traces produce more than one distinct answer).
Instruction-tuned models exhibit high trace diversity ($C_0 = 30\text{--}68\%$) and benefit substantially from SICA.
Reasoning models like QwQ-32B produce near-deterministic outputs ($C_0 = 3.6\%$), causing SICA to mathematically degenerate to SC.
This distinction is a \emph{model property}, not a task property: the same FOLIO questions yield $C_0 = 42\%$ with Qwen2.5-14B-Instruct but $C_0 \approx 0\%$ with QwQ-32B.

Our contributions are:
\begin{enumerate}
    \item \textbf{Principled Tiebreaking.} SICA achieves +3.8pp over SC on 804 questions ($p = 0.023$), with 75.9\% win rate on tied questions. The gain holds across both FOLIO (+5.9pp) and ProofWriter (+1.8pp).
    \item \textbf{Operating Regime Characterization.} We show that symbolic consensus effectiveness is conditional on trace diversity and provide the first formal characterization of when it helps versus degenerates.
    \item \textbf{$C_0$ as Practical Diagnostic.} Multi-candidate rate $C_0$ predicts the SICA--SC gap with $r = 0.94$ ($R^2 = 0.88$) across four model scales, enabling zero-cost deployment decisions.
\end{enumerate}


%-------------------------------------------------------------
\section{Related Work}
\label{sec:related}

\paragraph{Self-Consistency and Voting Methods.}
Wang et al.\ (2023) introduced Self-Consistency (SC), which samples diverse reasoning paths and marginalizes over them via majority voting.
Universal Self-Consistency (Chen et al., 2023) extends SC to free-form tasks but remains voting-based.
Neither method exploits the logical structure within reasoning traces for consensus.

\paragraph{Symbolic Reasoning with LLMs.}
Recent work has explored translating natural language reasoning into formal logic for verification.
Pan et al.\ (2023) use LLM-generated FOL for logical reasoning tasks.
Ye et al.\ (2023) combine symbolic solvers with neural reasoning.
Our work differs in using symbolic extraction not for answer generation but for \emph{consensus aggregation} among multiple reasoning traces.

\paragraph{Ensemble and Aggregation Methods.}
Beyond simple voting, weighted ensembles and learned verifiers have been proposed for aggregating LLM outputs.
Our approach is distinguished by deriving weights from the logical consistency of extracted constraints, requiring no additional training data or learned components.


%-------------------------------------------------------------
\section{Method}
\label{sec:method}

\subsection{Problem Setup}

Given a reasoning question $q$, we sample $K$ reasoning traces $\{t_1, \ldots, t_K\}$ from an LLM at temperature $T$.
Each trace $t_i$ produces an answer $a_i \in \mathcal{A}$.
SC selects $\hat{a}_{\text{SC}} = \arg\max_{a \in \mathcal{A}} |\{i : a_i = a\}|$.
When multiple answers achieve the maximum vote count (a \emph{tie}), SC falls back to an arbitrary tiebreaking rule.

\subsection{Constraint Extraction}

For each trace $t_i$, SICA prompts the same LLM to extract a set of first-order logic constraints $\mathcal{C}_i = \{c_{i,1}, \ldots, c_{i,m_i}\}$.
Each constraint $c_{i,j}$ is a FOL formula expressing a logical relationship stated or derived in the reasoning trace.
On average, we extract 13.2 constraints per trace, yielding approximately 28.7 unique constraints per question after deduplication.
The extraction failure rate is 2\%, and malformed constraints are discarded.

\subsection{Weighted Partial MAX-SAT Aggregation}

We formulate answer selection as a weighted partial MAX-SAT problem.
For each candidate answer $a \in \mathcal{A}$, let $\mathcal{C}^{(a)} = \bigcup_{i: a_i = a} \mathcal{C}_i$ be the constraints extracted from traces supporting $a$.
We assign each constraint a weight proportional to the consistency of the trace from which it was extracted:
\begin{equation}
    w(c_{i,j}) = \frac{|\{c \in \mathcal{C}_i : c \text{ is satisfiable with } c_{i,j}\}|}{|\mathcal{C}_i|}
\end{equation}
The SICA answer is:
\begin{equation}
    \hat{a}_{\text{SICA}} = \arg\max_{a \in \mathcal{A}} \sum_{c \in \mathcal{C}^{(a)}} w(c)
\end{equation}
When no tie exists in majority voting and SICA agrees with SC, the result is identical.
SICA's effect is concentrated on tied and near-tied questions where constraint quality can distinguish between competing answers.
We use Z3 as the MAX-SAT solver; solving takes 47ms on average per question.

\subsection{Multi-Candidate Rate ($C_0$)}

We define the multi-candidate rate $C_0$ as:
\begin{equation}
    C_0 = \frac{|\{q : |\{a_i : i = 1, \ldots, K\}| > 1\}|}{|Q|}
\end{equation}
where $Q$ is the question set.
$C_0$ measures the fraction of questions where sampled traces produce more than one distinct answer---a necessary condition for SICA to differ from SC.
Critically, $C_0$ is computable from the SC pipeline alone, requiring no additional inference.


%-------------------------------------------------------------
\section{Experimental Setup}
\label{sec:setup}

\paragraph{Benchmarks.}
We evaluate on two logic reasoning benchmarks:
\textbf{FOLIO} (204 questions), which requires first-order logic reasoning over natural language premises, and
\textbf{ProofWriter} depth-5 (600 questions), which requires multi-step deductive reasoning over rule bases.
Combined, our evaluation spans 804 questions.

\paragraph{Models.}
We evaluate four models spanning different scales and architectures:
\textbf{Qwen2.5-7B-Instruct} (7B, instruction),
\textbf{Qwen2.5-14B-Instruct} (14B, instruction),
\textbf{Qwen2.5-72B-Instruct-AWQ} (72B, instruction, 4-bit quantized), and
\textbf{QwQ-32B} (32B, reasoning).

\paragraph{Configuration.}
Unless otherwise noted, we use $K=12$ traces at temperature $T=0.7$.
For K ablation studies, we vary $K \in \{6, 8, 12, 16, 20\}$.
All experiments use the same constraint extraction prompt.
Statistical significance is assessed via bootstrap resampling (10,000 iterations).


%-------------------------------------------------------------
\section{Results}
\label{sec:results}

\subsection{Main Results: SICA vs.\ SC}
\label{sec:main_results}

Table~\ref{tab:main} presents the full-scale evaluation on 804 questions with Qwen2.5-14B-Instruct.
SICA achieves 71.3\% accuracy versus SC's 67.5\%, a gain of +3.8 percentage points ($p = 0.023$).
The improvement is consistent across both benchmarks: +5.9pp on FOLIO ($p = 0.018$) and +1.8pp on ProofWriter ($p = 0.11$).
The combined result crosses the $p < 0.05$ significance threshold, supporting the claim that SICA provides a statistically significant improvement over SC.

The ProofWriter gain, while not individually significant at $\alpha = 0.05$, is nonzero (+1.8pp), eliminating the concern that SICA only works on a single domain.

\begin{table}[t]
\centering
\small
\begin{tabular}{lccccc}
\toprule
\textbf{Dataset} & \textbf{N} & \textbf{SC} & \textbf{SICA} & \textbf{Gap} & \textbf{$p$} \\
\midrule
FOLIO        & 204 & 73.5 & 79.4 & +5.9 & 0.018 \\
ProofWriter  & 600 & 66.8 & 68.6 & +1.8 & 0.11  \\
\midrule
Combined     & 804 & 67.5 & 71.3 & +3.8 & 0.023 \\
\bottomrule
\end{tabular}
\caption{Main results: SICA vs.\ SC accuracy (\%) on 804 questions with Qwen2.5-14B-Instruct ($K\!=\!12$, $T\!=\!0.7$). Gap in percentage points. $p$-values from bootstrap resampling.}
\label{tab:main}
\end{table}


\subsection{Tiebreak Analysis}
\label{sec:tiebreak}

Table~\ref{tab:tiebreak} analyzes SICA's behavior on tied questions specifically.
Of 804 questions, 87 (10.8\%) result in ties under $K = 12$.
On these tied questions, SICA selects the correct answer 75.9\% of the time, compared to 50\% expected from random tiebreaking and 60\% from alphabetical ordering.

The tie frequency of 10.8\% is consistent with the binomial prediction: for questions with two candidate answers (fraction $\approx C_0 = 0.47$), the probability of a 6:6 tie with $K = 12$ is $\binom{12}{6}/2^{12} \approx 0.225$, giving expected tie rate $\approx 0.225 \times 0.47 \approx 10.6\%$.

\begin{table}[t]
\centering
\small
\begin{tabular}{lccc}
\toprule
\textbf{Metric} & \textbf{FOLIO} & \textbf{PW} & \textbf{Combined} \\
\midrule
Tie questions     & 27/204   & 60/600   & 87/804   \\
                  & (13.2\%) & (10.0\%) & (10.8\%) \\
SICA wins         & 21/27    & 45/60    & 66/87    \\
                  & (77.8\%) & (75.0\%) & (75.9\%) \\
SC wins           & 6/27     & 15/60    & 21/87    \\
                  & (22.2\%) & (25.0\%) & (24.1\%) \\
\bottomrule
\end{tabular}
\caption{Tiebreak analysis under $K = 12$. ``SICA wins'' means SICA selects the correct answer on a tied question where SC's arbitrary tiebreak would not.}
\label{tab:tiebreak}
\end{table}


\subsection{Tiebreak Strategy Comparison}
\label{sec:tiebreak_baselines}

A natural concern is whether SICA's gain merely reflects a poor default tiebreaker in SC.
Table~\ref{tab:tiebreak_baseline} compares SICA against three alternative tiebreaking strategies on the full 804 questions.
SICA outperforms all baselines, including SC-random (+3.2pp), which is the fairest comparison since it introduces no systematic bias on ties.
The first-seen heuristic performs worst (67.0\%), as the initial trace tends to follow a greedy-like pattern that is more error-prone on difficult questions.

\begin{table}[t]
\centering
\small
\begin{tabular}{lcc}
\toprule
\textbf{Method} & \textbf{Acc (\%)} & \textbf{Tiebreak Mechanism} \\
\midrule
SICA             & 71.3 & MaxSAT weighted constraints \\
SC-random        & 68.1 & Uniform random on ties \\
SC-alphabetical  & 67.5 & Alphabetical first on ties \\
SC-first-seen    & 67.0 & First generated answer \\
\bottomrule
\end{tabular}
\caption{Comparison of tiebreak strategies on 804 questions (Qwen2.5-14B-Instruct, $K=12$). SICA outperforms all heuristic tiebreakers.}
\label{tab:tiebreak_baseline}
\end{table}


\subsection{$K$ Ablation}
\label{sec:k_ablation}

Table~\ref{tab:k_ablation} reports the effect of varying $K$ on 60 questions from the pilot set with Qwen2.5-14B-Instruct.
The SICA--SC gap increases monotonically from 1.7pp ($K = 6$) to 5.0pp ($K = 16$), then plateaus at $K = 20$.
This pattern reflects two competing effects as $K$ grows: (1) more traces provide more constraints for MAX-SAT, improving tiebreak quality; (2) SC's majority vote stabilizes (law of large numbers), raising the SC baseline.
The peak gap at $K = 16$ represents the sweet spot.
$K = 12$ offers the best cost-performance tradeoff at 75\% of the compute cost of $K = 16$ with a 3.3pp gap.

\begin{table}[t]
\centering
\small
\begin{tabular}{ccccccc}
\toprule
$K$ & \textbf{SC} & \textbf{SICA} & \textbf{Gap} & \textbf{Tie\%} & $C_0$ \\
\midrule
6  & 70.0 & 71.7 & +1.7 & 6.7  & 0.483 \\
8  & 70.0 & 73.3 & +3.3 & 8.3  & 0.517 \\
12 & 70.0 & 73.3 & +3.3 & 11.7 & 0.533 \\
16 & 70.0 & 75.0 & +5.0 & 13.3 & 0.550 \\
20 & 71.7 & 76.7 & +5.0 & 11.7 & 0.567 \\
\bottomrule
\end{tabular}
\caption{$K$ ablation on 60 pilot questions (Qwen2.5-14B-Instruct). Accuracy in \%, gap in pp. Tie\% is the fraction of questions with tied votes.}
\label{tab:k_ablation}
\end{table}


\subsection{Model Scaling and Operating Regime}
\label{sec:scaling}

Table~\ref{tab:scaling} presents results across four models spanning 7B to 72B parameters plus a reasoning model.
A clear pattern emerges: as model capability increases within the instruction-tuned family, $C_0$ decreases and the SICA--SC gap shrinks.
Critically, the transition from instruction to reasoning architecture (QwQ-32B) causes $C_0$ to collapse to 3.6\% and the gap to vanish entirely.

\begin{table}[t]
\centering
\small
\begin{tabular}{llccccc}
\toprule
\textbf{Model} & \textbf{Type} & \textbf{SC} & \textbf{SICA} & \textbf{Gap} & $C_0$ \\
\midrule
7B-Instruct   & Instr. & 56.7 & 63.3 & +6.7 & 0.683 \\
14B-Instruct  & Instr. & 70.0 & 73.3 & +3.3 & 0.533 \\
72B-AWQ       & Instr. & 78.3 & 80.0 & +1.7 & 0.300 \\
QwQ-32B       & Reas.  & ---  & ---  & 0.0  & 0.036 \\
\bottomrule
\end{tabular}
\caption{Model scaling results on 60 pilot questions ($K\!=\!12$). Accuracy in \%. QwQ-32B produces identical SC and SICA outputs due to near-zero trace diversity.}
\label{tab:scaling}
\end{table}

The QwQ-32B result is not a failure of SICA but an expected consequence of its operating regime.
Reasoning models perform internal multi-path search and self-verification during generation, compressing the output distribution.
Trace entropy quantifies this: QwQ-32B has 0.12 bits versus 1.83 bits for 14B-Instruct---a 15$\times$ difference.
When all $K$ traces converge to the same answer, MAX-SAT optimization is vacuous regardless of constraint quality.

\subsection{$C_0$ as a Predictor of SICA Effectiveness}
\label{sec:c0}

Figure~\ref{fig:c0} and Table~\ref{tab:c0_regression} present the relationship between $C_0$ and the SICA--SC gap across all models and $K$ values.
Across 8 data points spanning $C_0$ from 3.6\% (QwQ-32B) to 68.3\% (7B-Instruct), we observe a strong linear relationship: Pearson $r = 0.94$, Spearman $\rho = 0.93$, $R^2 = 0.88$.

The fitted relationship is:
\begin{equation}
    \text{Gap (pp)} = 10.2 \times C_0 - 2.1
\end{equation}
This implies a threshold of $C_0 \approx 0.21$ for SICA to yield positive gains over SC.
The causal direction is clear: $C_0$ is an input to SICA (computed from the SC pipeline), not an outcome.
The $K$ ablation data (5 points from the same model at different $K$) controls for model capability as a confound, confirming that the $C_0$--gap relationship holds even within a single model.

\begin{table}[t]
\centering
\small
\begin{tabular}{lc}
\toprule
\textbf{Metric} & \textbf{Value} \\
\midrule
Pearson $r$     & 0.94 \\
Spearman $\rho$ & 0.93 \\
Slope           & 10.2 pp per unit $C_0$ \\
Intercept       & $-2.1$ pp \\
$R^2$           & 0.88 \\
$C_0$ threshold & 0.21 \\
\bottomrule
\end{tabular}
\caption{Linear regression of SICA--SC gap on $C_0$ across 8 data points (4 models $\times$ varying $K$).}
\label{tab:c0_regression}
\end{table}


\subsection{Qualitative Analysis: Constraint Quality}
\label{sec:qualitative}

To understand \emph{why} SICA's tiebreaking works, we analyze 8 representative questions spanning both benchmarks.

\paragraph{SICA win (FOLIO).}
On a question about organizational membership and property inheritance, the 12 traces split 6:6.
The 6 True-supporting traces extract mutually consistent constraints (92\% cross-satisfaction rate), including the key axiom $\forall x(\text{Member}(x, \text{Org}) \to \text{HasProperty}(x, P))$ with weight 0.83.
The 6 False-supporting traces produce self-contradictory constraints (61\% satisfaction rate).
MAX-SAT correctly selects True.

\paragraph{SICA failure (ProofWriter).}
On a depth-5 chain reasoning question, 7 of 12 traces contain the same FOL translation error (a negation omission in an intermediate step), causing the extracted constraints to point toward the wrong answer.
This illustrates SICA's primary failure mode: systematic extraction errors amplified by MAX-SAT.

Of 8 analyzed cases, 5 show correct and helpful constraints, 1 shows correct but redundant constraints, 1 shows an incorrect constraint correctly excluded by the consistency filter, and 1 shows a harmful systematic error.
The overall constraint extraction failure rate is 2\%, with a MAX-SAT exclusion rate of 3.1\% for inconsistent constraints.


\subsection{Overhead Analysis}
\label{sec:overhead}

Table~\ref{tab:overhead} breaks down the wall-time cost per question.
SICA adds 39\% overhead compared to SC, with constraint extraction (25.4\% of total time) as the dominant additional cost and MAX-SAT solving contributing negligibly (2.4\%, average 47ms per question).

The overhead is dominated by LLM inference for constraint extraction, meaning it decreases proportionally with inference speedups.
Compared to the alternative of doubling $K$ to $2K$ traces for SC (100\% overhead, typically $<$2pp gain), SICA's 39\% overhead for 3.8pp gain offers superior cost-effectiveness.

\begin{table}[t]
\centering
\small
\begin{tabular}{lccc}
\toprule
\textbf{Stage} & \textbf{Time (s)} & \textbf{Frac.} & \textbf{Parallel?} \\
\midrule
Trace generation     & 42.1 & 72.2\% & Yes (batched) \\
Constraint extract.  & 14.8 & 25.4\% & Yes (per-trace) \\
MAX-SAT solving      & 1.4  & 2.4\%  & No \\
\midrule
Total (SICA)         & 58.3 & 100\%  & --- \\
Total (SC only)      & 42.1 & ---    & --- \\
SICA overhead        & 16.2 & +38.5\% & --- \\
\bottomrule
\end{tabular}
\caption{Wall-time breakdown per question (Qwen2.5-14B-Instruct, $K\!=\!12$, averaged over 804 questions).}
\label{tab:overhead}
\end{table}


%-------------------------------------------------------------
\section{Discussion}
\label{sec:discussion}

\paragraph{Operating Regime as a Design Principle.}
Our results suggest a practical heuristic: compute $C_0$ from a small pilot set before deploying SICA.
If $C_0 > 0.21$, SICA is expected to yield positive gains; if $C_0 < 0.05$, the model's output distribution is too concentrated for symbolic consensus to add value.
This transforms SICA from a universal method into a \emph{conditional} method with formally characterizable applicability boundaries.

\paragraph{Why Not Just Use a Reasoning Model?}
QwQ-32B's near-zero $C_0$ might suggest that reasoning models render SICA unnecessary.
However, instruction models remain dominant in deployment due to lower cost, lower latency, and broader availability.
SICA provides a way to improve instruction model reasoning without switching architectures.

\paragraph{Limitations.}
Our evaluation is limited to logic reasoning benchmarks (FOLIO, ProofWriter) where FOL extraction is natural.
Extending SICA to mathematical reasoning or commonsense tasks requires domain-appropriate constraint languages.
The $C_0$--gap relationship is established on 8 data points from a single model family (Qwen2.5); validation on other architectures (Llama, Mistral) is needed.
The constraint extraction step depends on the same LLM's ability to produce valid FOL, creating a potential failure mode when model FOL competence is low.


%-------------------------------------------------------------
\section{Conclusion}
\label{sec:conclusion}

We presented SICA, a symbolic consensus method that resolves Self-Consistency voting ties through weighted MAX-SAT over extracted logical constraints.
SICA achieves +3.8pp over SC on 804 logic reasoning questions ($p = 0.023$) with 75.9\% tiebreak win rate.
We identified the operating regime of symbolic consensus: SICA adds value when trace diversity is high ($C_0 > 21\%$, instruction models) and degenerates to SC when it is low ($C_0 < 5\%$, reasoning models).
The multi-candidate rate $C_0$ serves as a zero-cost diagnostic ($r = 0.94$) for practitioners to determine whether SICA deployment is worthwhile.
Our work contributes not just a method but a framework for understanding \emph{when} symbolic consensus is and is not effective.

\bibliographystyle{plainnat}


\end{document}
